\begin{abstract}
Mathematical reasoning remains a formidable obstacle for language models due to its inherent complexity and highly structured requirements. In this work, we present DeepSeekMath~7B, a model developed by further pre-training DeepSeek-Coder-Base-v1.5 7B on 120B mathematics-focused tokens gathered from Common Crawl, complemented by both natural language and programming code sources. Without leveraging external mathematical toolkits or ensemble methods, DeepSeekMath~7B attains a strong 51.7\% score on the MATH benchmark at the competition level, narrowing the gap with Gemini-Ultra and GPT-4. When utilizing self-consistency with 64 generated samples, DeepSeekMath~7B achieves a result of 60.9\% on the same benchmark. The advanced mathematical reasoning abilities of DeepSeekMath~stem from two principal elements: (1) the extensive exploitation of publicly available online data by means of a carefully designed data selection procedure, and (2) the adoption of Group Relative Policy Optimization (GRPO)—a new adaptation of Proximal Policy Optimization (PPO)—which simultaneously bolsters mathematical reasoning proficiency and optimizes PPO's memory efficiency.
\end{abstract}

\section{Introduction}

Large language models (LLM) have drastically transformed how artificial intelligence tackles mathematical reasoning tasks, yielding remarkable improvements on benchmarks assessing quantitative reasoning~\citep{MATH} and geometry reasoning~\citep{trinh2024solving}. In addition, these models increasingly serve as valuable aids to humans confronting challenging mathematical questions~\citep{tao}. Nevertheless, leading systems such as GPT-4~\citep{gpt4} and Gemini-Ultra~\citep{gemini} remain closed-source, and open-source alternatives continue to lag substantially in their mathematical performance.

This work presents DeepSeekMath, a specialized language model tailored for mathematical tasks, which demonstrates superior mathematical performance relative to existing open-source models and closely matches GPT-4 on standardized academic evaluations. The model's capabilities are enabled by the development of the DeepSeekMath Corpus, an extensive, high-quality pre-training dataset composed of 120B mathematical tokens. This corpus is compiled from Common Crawl (CC) data utilizing a classifier built with fastText \citep{joulin2016fasttext}. For the initial phase, the classifier is trained with positive examples from OpenWebMath \citep{openwebmath} and negative examples sourced from various unrelated web content. Using this classifier, further positive samples are identified within CC and subsequently curated through human review. The improved dataset is used to retrain the classifier, thereby enhancing its effectiveness. Empirical evaluations demonstrate that the resulting dataset supports high-caliber model performance: the DeepSeekMath-Base 7B model achieves 64.2\% accuracy on GSM8K \citep{gsm8k} and 36.2\% on the MATH competition benchmark \citep{MATH}, exceeding the results of Minerva 540B \citep{minerva}. Furthermore, due to the multilingual nature of the DeepSeekMath Corpus, advancements are observed on Chinese mathematics benchmarks \citep{wei2023cmath,agieval}. We propose that our methods for mathematical data curation offer a valuable foundation for subsequent research efforts, with substantial opportunities for further advancement.

DeepSeekMath-Base is built upon DeepSeek-Coder-Base-v1.5 7B \citep{deepseek-coder}, as our findings indicate that initializing with a model trained on code yields superior results compared to starting from a general-purpose LLM. Additionally, we find that mathematical pre-training not only augments mathematical proficiency but also enhances the model’s performance on general benchmarks such as MMLU \citep{mmlu} and BBH \citep{bbh}, reflecting broader gains in reasoning capability.

Subsequent to pre-training, we subject DeepSeekMath-Base to a phase of mathematical instruction tuning using chain-of-thought \citep{cot}, program-of-thought \citep{pot,pal}, and tool-augmented reasoning datasets \citep{tora}. The fine-tuned model, DeepSeekMath-Instruct 7B, surpasses all comparable 7B-sized models and demonstrates performance on par with instruction-tuned models of the 70B scale.

We also present Group Relative Policy Optimization (GRPO), an adaptation of the Proximal Policy Optimization (PPO) algorithm \citep{schulman2017proximal}, for reinforcement learning (RL). GRPO eliminates the need for a critic by inferring the baseline from a collection of grouped scores, leading to notable reductions in computational demands during training. Leveraging only a portion of the English instruction-tuning corpus, GRPO delivers pronounced improvements over DeepSeekMath-Instruct, as reflected in both in-domain tasks (GSM8K: 82.9\% $\rightarrow$ 88.2\%, MATH: 46.8\% $\rightarrow$ 51.7\%) and out-of-domain mathematical challenges (e.g., CMATH: 84.6\% $\rightarrow$ 88.8\%) during the RL stage. Additionally, we develop a cohesive framework to compare and interpret diverse methods, such as Rejection Sampling Fine-Tuning (RFT) \citep{yuan2023scaling}, Direct Preference Optimization (DPO) \citep{dpo}, PPO, and GRPO. Our unified perspective shows that all these techniques can be viewed as either direct or streamlined RL methodologies. We also perform comprehensive evaluations—including contrasts between online and offline training, outcome versus process supervision, and single-turn versus iterative RL—to systematically examine the core components of this unified approach. Finally, we discuss the mechanisms through which our RL strategy enhances instruction-tuned model performance, and outline future research directions for developing more efficient RL strategies under this unified paradigm.

\subsection{Contributions}
This work presents advancements in scalable math pre-training, together with an investigation and assessment of reinforcement learning methods.

\textbf{Math Pre-Training at Scale}

\begin{itemize}[topsep=0pt]
    \item Our study demonstrates that publicly available Common Crawl data possesses significant utility for mathematical applications. Through a carefully engineered data filtering pipeline, we assemble the DeepSeekMath Corpus—a rigorously curated collection containing 120B tokens from mathematically relevant web content. This corpus surpasses the math web page datasets of Minerva \citep{minerva} by nearly a factor of seven, and the recently introduced OpenWebMath \citep{openwebmath} by a factor of nine in terms of size.

    \item The pre-trained DeepSeekMath-Base 7B model attains performance on par with Minerva 540B \citep{minerva}, demonstrating that the scale of model parameters is not the sole determinant for excelling at mathematical reasoning. Our results indicate that smaller models, when pre-trained on high-caliber data, can also yield robust mathematical capabilities.

    \item We provide insights derived from our mathematical training experiments, revealing that performing code training prior to math training boosts models’ effectiveness at solving math problems, irrespective of tool usage. This finding contributes to the ongoing debate over whether code pre-training enhances reasoning proficiency; we posit that, at least for mathematics, it indeed does.

    \item In spite of the prevalent use of arXiv papers for training—particularly in domains centered on mathematics—our experiments show that such training offers no substantial benefit on the array of mathematical benchmarks evaluated in this study.
\end{itemize}

\textbf{Exploration and Analysis of Reinforcement Learning}

\begin{itemize}[topsep=0pt]
    \item We present Group Relative Policy Optimization (GRPO), a novel and resource-efficient reinforcement learning algorithm. GRPO eliminates the need for a critic network by leveraging group-based score baselines, which leads to substantial reductions in training overhead when compared to Proximal Policy Optimization (PPO).
    \item Our results show that GRPO markedly improves the capabilities of our instruction-tuned model, DeepSeekMath-Instruct, utilizing only the instruction-tuning dataset. Additionally, we detect gains in the model's generalization to out-of-domain tasks during the reinforcement learning phase.
    \item We offer a cohesive framework that elucidates the relationships among various methods, including RFT, DPO, PPO, and GRPO. Through comprehensive experimentation—such as comparing online and offline training strategies, evaluating outcome versus process supervision, and analyzing single-turn versus iterative reinforcement learning—we systematically examine the key factors underpinning this paradigm.
    \item Building upon our unified framework, we investigate the underlying mechanisms that contribute to the effectiveness of reinforcement learning and outline several promising avenues for enhancing reinforcement learning approaches for large language models.
\end{itemize}

\subsection{Summary of Evaluations and Metrics}

\begin{itemize}[topsep=0pt]
    \item \textbf{English and Chinese Mathematical Reasoning}: Our evaluation encompasses a diverse range of benchmarks for both English and Chinese, spanning mathematics problems from elementary to undergraduate level. The English benchmarks consist of GSM8K \citep{gsm8k}, MATH \citep{MATH}, SAT \citep{llemma}, OCW Courses \citep{minerva}, and MMLU-STEM \citep{mmlu}. The Chinese benchmarks feature MGSM-zh \citep{mgsm}, CMATH \citep{wei2023cmath}, Gaokao-MathCloze \citep{agieval}, and Gaokao-MathQA \citep{agieval}. We assess the models’ proficiency in independently generating complete textual solutions as well as their problem-solving skills utilizing Python.

    For the English datasets, DeepSeekMath-Base achieves competitive performance relative to the proprietary Minerva 540B \citep{minerva}, and it outperforms all currently available open-source base models—such as Mistral 7B \citep{mistral} and Llemma-34B \citep{llemma}—by a considerable margin, irrespective of previous math-focused pre-training. DeepSeekMath-Base particularly excels on Chinese evaluations, which can be attributed to our inclusion of high-quality multilingual mathematical data, diverging from the English-only pre-training approach of prior studies \citep{minerva,llemma}. Through mathematical instruction tuning and reinforcement learning, the upgraded models, DeepSeekMath-Instruct and DeepSeekMath-RL, deliver robust results, becoming the first open-source models to exceed the 50\% accuracy threshold on the challenging competition-level MATH benchmark.
\end{itemize}

\item \textbf{Formal Mathematics}: For the assessment of DeepSeekMath-Base, we adopt the informal-to-formal theorem proving benchmark introduced by \citep{dsp_proof}, employing the miniF2F dataset \citep{minif2f} and selecting Isabelle \citep{isabelle} as the proof assistant. Our results indicate that DeepSeekMath-Base achieves notable few-shot performance in autoformalization tasks.

\item \textbf{Natural Language Understanding, Reasoning, and Code}: To obtain a comprehensive evaluation of general comprehension, logical reasoning, and programming skills, we test DeepSeekMath-Base using the Massive Multitask Language Understanding (MMLU) suite \citep{mmlu}, which includes 57 multiple-choice questions from various domains; BIG-Bench Hard (BBH) \citep{bbh}, featuring 23 complex problems that typically necessitate multi-step reasoning; and well-established programming benchmarks HumanEval \citep{codex} and MBPP \citep{mbpp}. The results confirm that pre-training on mathematical data significantly enhances both reasoning and language understanding capacities.

\section{Math Pre-Training}

\subsection{Data Collection and Decontamination} This section details the methodology behind assembling the DeepSeekMath Corpus using Common Crawl data. As illustrated in Figure \ref{fig:data_collect}, we introduce an iterative approach that allows for systematic extraction of an extensive mathematical text corpus, beginning with a foundational seed dataset (e.g., a compact and curated collection of mathematics-related texts). Notably, this methodology is adaptable to additional domains, such as software engineering.

Initially, OpenWebMath \citep{openwebmath}, a rigorously filtered compilation of mathematical web content, serves as our seed dataset. Leveraging this collection, we develop a fastText model \citep{joulin2016fasttext} to retrieve additional web documents exhibiting properties akin to those in OpenWebMath. We sample 500,000 entries at random from the seed as positive examples, pairing them with 500,000 randomly selected Common Crawl webpages designated as negative samples. Model training utilizes a public implementation\footnote{\url{https://fasttext.cc}}, with the embedding dimension set to 256, a learning rate of 0.1, a maximum word n-gram length of 3, a minimum word count of 3, and training over 3 epochs. To further constrain the volume of Common Crawl, we utilize both URL-level deduplication and near-duplicate elimination, thus reducing the dataset to 40B HTML pages. Next, the fastText model is used to retrieve relevant mathematical content from this deduplicated web data. To eliminate subpar mathematical material, gathered pages are ranked by fastText model predictions, retaining only those with the highest scores. The chosen dataset volume is empirically validated through pre-training trials on sets of the top 40B, 80B, 120B, and 160B tokens. For the initial pass, we select the top 40B tokens for further use.

Following the initial round of data gathering, a significant number of mathematical web pages remain uncollected. This shortfall can primarily be attributed to the fastText model being trained on a collection of positive samples that does not sufficiently represent the diversity of mathematical web content. To remedy this, we augment the seed corpus by incorporating a broader array of mathematical web sources, thus enhancing the fastText model's coverage. Our strategy begins with partitioning the complete Common Crawl dataset into non-overlapping domains, where each domain comprises web pages sharing an identical base URL. For every domain, we compute the proportion of web pages obtained during the first iteration. Domains from which over 10\% of web pages have been retrieved are categorized as math-centric (for example, \url{mathoverflow.net}). Next, within these domains, we manually identify URLs that correspond to mathematical resources (such as \url{mathoverflow.net/questions}). Any previously uncollected web pages linked to these mathematical URLs are added to the seed corpus. Through this method, we acquire a larger and more varied set of positive examples, thereby allowing the fastText model to better identify mathematical content in future iterations. After four collection cycles, this process yields 35.5M mathematical web pages, amounting to a total of 120B tokens. Notably, in the fourth round, we observe that nearly 98\% of the collected data overlaps with the third iteration, prompting us to terminate further data collection at that stage.

To prevent contamination of benchmark datasets, we adopt the filtering strategy described in \cite{deepseek-coder}. Specifically, we remove web pages that include questions or answers taken from English mathematics benchmarks such as GSM8K \citep{gsm8k} and MATH \citep{MATH}, as well as Chinese benchmarks like CMATH \citep{wei2023cmath} and AGIEval \citep{agieval}. The applied filtering mechanism is as follows: any segment of text that contains an exact match to a 10-gram sequence from these evaluation benchmarks is excluded from our mathematical training data. For benchmark sequences with fewer than 10 but at least 3 grams, exact string matching is similarly used to eliminate contaminated content.

\subsection{Assessing the DeepSeekMath~Corpus Quality}

We conduct pre-training studies to compare the DeepSeekMath~Corpus against several recently introduced mathematical training corpora:

\begin{itemize}[topsep=0pt]
    \item \textbf{MathPile} \citep{mathpile}: This 8.9B-token corpus aggregates content from multiple sources including textbooks, Wikipedia, ProofWiki, CommonCrawl, StackExchange, and arXiv, with over 85\% of data coming from arXiv;
    \item \textbf{OpenWebMath} \citep{openwebmath}: Derived by filtering CommonCrawl for mathematical material, resulting in a dataset with 13.6B tokens;
    \item \textbf{Proof-Pile-2} \citep{llemma}: This mathematical dataset brings together OpenWebMath, AlgebraicStack (10.3B tokens of math code), and arXiv papers (28.0B tokens). For experiments using Proof-Pile-2, we adhere to the arXiv:Web:Code ratio of 2:4:1 as described in \cite{llemma}.
\end{itemize}

\subsubsection{Training Setting}
\label{sec:quality-policy}
We conduct math-specific training on a general-purpose language model comprising 1.3 billion parameters, structured identically to the DeepSeek LLMs \citep{deepseek-llm}, which we refer to as DeepSeek-LLM 1.3B. For each mathematical dataset, we independently train a separate model for 150 billion tokens. All training procedures utilize the efficient HAI-LLM \citep{haillm} training framework. Consistent with DeepSeek LLMs’ established training regimen, we adopt the AdamW optimizer \citep{adamW}, setting $\beta_1=0.9$, $\beta_2=0.95$, and $\mathrm{weight\_decay}=0.1$. The learning rate follows a multi-phase schedule: it linearly increases during 2{,}000 warmup steps to its maximum, then decays to 31.6\% of the peak value after 80\% of the training steps, and further decreases to 10.0\% at 90\% of the training completion. The maximum learning rate is configured as 5.3e-4, with training carried out using a batch size of 4 million tokens and a context window of 4,000 tokens.

\subsubsection{Evaluation Results}

\textbf{The DeepSeekMath~Corpus stands out due to its exceptional quality, comprehensive multilingual coverage, and unparalleled scale.}

\begin{itemize}[topsep=0pt]
    \item \textbf{High-quality}:
    Downstream evaluation is performed across eight mathematical tasks using few-shot chain-of-thought prompting \cite{cot}. According to Table \ref{tab:corpora_comparison}, models fine-tuned on DeepSeekMath~Corpus consistently surpass others in performance. Moreover, as illustrated in Figure \ref{fig:corpora_comparisons}, the DeepSeekMath Corpus-trained model already outperforms those trained on Proof-Pile-2 after 50B tokens (corresponding to a complete epoch over Proof-Pile-2), implying that DeepSeekMath Corpus exhibits a higher average data quality.
    \item \textbf{Multilingual}:
    DeepSeekMath~Corpus incorporates mathematical text spanning various languages, with English and Chinese making up the largest fractions. Table \ref{tab:corpora_comparison} reveals that pre-training with DeepSeekMath~Corpus significantly improves mathematical reasoning in both English and Chinese. In contrast, other mathematical datasets, being predominantly English-only, deliver minimal benefit to, or may even impede, Chinese mathematical reasoning.
    \item \textbf{Large-scale}:
    Relative to existing mathematical corpora, DeepSeekMath~Corpus offers a considerably greater volume of data. As seen in Figure \ref{fig:corpora_comparisons}, DeepSeek-LLM 1.3B exhibits a faster learning trajectory and sustains progress longer when trained on DeepSeekMath~Corpus. Baseline datasets, being much more limited in size, are exposed repeatedly during training, causing their corresponding models to plateau rapidly.
\end{itemize}

\subsection{Training and Evaluating DeepSeekMath-Base 7B}

This section details DeepSeekMath-Base 7B, a foundational model designed for advanced reasoning, with a particular emphasis on mathematical competence. The model's initialization leverages DeepSeek-Coder-Base-v1.5 7B \citep{deepseek-coder}, and it undergoes training for a total of 500B tokens. The composition of the training dataset consists of 56\% from the DeepSeekMath Corpus, 4\% contributed by AlgebraicStack, 10\% sourced from arXiv, 20\% derived from code repositories on Github, and the final 10\% comprising natural language data in both English and Chinese extracted from Common Crawl. The training procedure predominantly follows the setup in Section \ref{sec:quality-policy}, with modifications including a maximum learning rate set to 4.2e-4 and a batch size configured to 10M tokens.

We systematically evaluate DeepSeekMath-Base 7B's proficiency in mathematical reasoning, considering its capacity to independently generate complete solutions, utilize auxiliary tools for problem-solving, and perform rigorous formal theorem proving. In addition to its mathematical specialization, we also report the model's broader abilities, assessing competencies in general language understanding, logical reasoning, and programming.

\paragraph{Mathematical Problem Solving with Step-by-Step Reasoning}

To assess DeepSeekMath-Base's problem-solving prowess, we conduct few-shot chain-of-thought prompting \citep{cot} experiments on eight benchmarks, spanning both English and Chinese. These benchmarks address a spectrum of mathematical reasoning tasks, such as quantitative reasoning (e.g., GSM8K \citep{gsm8k}, MATH \citep{MATH}, CMATH \citep{wei2023cmath}) and multiple-choice formats (e.g., MMLU-STEM \citep{mmlu}, Gaokao-MathQA \citep{agieval}), and they span various mathematical disciplines ranging from elementary to undergraduate-level topics.

As indicated in Table \ref{tab:base_math_cot}, DeepSeekMath-Base 7B achieves the highest performance across all eight benchmarks among open-source base models, including both the popular general-purpose Mistral 7B \citep{mistral} and the recently introduced, math-trained Llemma 34B \citep{llemma} (which leveraged Proof-Pile-2 \citep{llemma}). In particular, on the challenging competition-grade MATH benchmark, DeepSeekMath-Base attains an improvement of more than 10\% absolute compared to all prior open-source base models, and even exceeds the performance of the much larger (77x) closed-source Minerva 540B \citep{minerva}, which is an extension of PaLM \citep{palm} with additional mathematical text training.

\paragraph{Mathematical Problem Solving with Tool Use} For assessments involving program-aided mathematical reasoning, we evaluate on GSM8K and MATH by employing few-shot program-of-thought prompting strategies \citep{pot,pal}. In this setup, the model addresses each problem by authoring a Python script, drawing upon libraries such as \textit{math} and \textit{sympy} for handling complex computations. The resultant program output serves as the model’s answer. Table \ref{tab:base_math_pot_proof} demonstrates that DeepSeekMath-Base 7B surpasses the previous best open-source model, Llemma 34B, in this context.

\paragraph{Formal Mathematics}

Automating formal proofs has grown in relevance for bolstering both the correctness and efficiency of mathematical proofs. We assess DeepSeekMath-Base 7B on the informal-to-formal proof task from \citep{dsp_proof}, which involves generating a formal proof given an informal description, a formalized statement, and an informal proof sketch. The evaluation is conducted on miniF2F \citep{minif2f}, a dataset covering Olympiad-level formal mathematics, with the model producing Isabelle formalizations using few-shot prompts. Consistent with \cite{dsp_proof}, our pipeline has the model draft proof outlines, with Sledgehammer \citep{sledgehammer} subsequently employed to automate missing inference steps. Results in Table \ref{tab:base_math_pot_proof} highlight the strong autoformalization capabilities of DeepSeekMath-Base 7B.

\paragraph{Natural Language Understanding, Reasoning, and Code}

The model is further assessed for natural language understanding using MMLU \citep{mmlu}, reasoning proficiency via BBH \citep{bbh}, and programming ability on HumanEval \citep{codex} as well as MBPP \citep{mbpp}. Table \ref{tab:understanding_reasoning_code} shows that DeepSeekMath-Base 7B delivers marked improvements on MMLU and BBH compared to its predecessor, DeepSeek-Coder-Base-v1.5 \citep{deepseek-coder}, thus reflecting the advantages conferred by targeted mathematical training. Moreover, continual training with code tokens allows DeepSeekMath-Base 7B to match the coding benchmark performance of DeepSeek-Coder-Base-v1.5. Taken together, DeepSeekMath-Base 7B considerably outperforms Mistral 7B \citep{mistral} across all three evaluated reasoning and coding tasks.

\section{Supervised Fine-Tuning}

\subsection{SFT Data Curation}

We develop a dataset specifically for mathematical instruction tuning, comprising problems in both English and Chinese across a spectrum of mathematical domains and difficulty. Each problem is coupled with a solution articulated in formats such as chain-of-thought (CoT) \citep{cot}, program-of-thought (PoT) \citep{pot,pal}, or in tool-integrated reasoning style \citep{tora}. The curated set includes a total of 776,000 training instances.

\begin{itemize}[topsep=0pt]
    \item \textbf{English mathematical datasets}: For English-language data, we provide tool-integrated solution annotations for GSM8K and MATH, and also incorporate selected problems from MathInstruct \citep{MathInstruct} as well as the training partition of Lila-OOD \citep{lila}, where solutions are presented through CoT or PoT methodologies. The English datasets encompass a broad range of mathematical disciplines, such as algebra, probability, number theory, calculus, and geometry.
    \item \textbf{Chinese mathematical datasets}: Our Chinese corpus includes K-12 mathematics problems distributed across 76 subfields (for example, linear equations), with accompanying solutions documented in both CoT and tool-integrated formats.
\end{itemize}

\subsection{Training and Evaluating DeepSeekMath-Instruct 7B}

Here, we detail the instruction tuning process of DeepSeekMath-Instruct 7B, which is trained atop DeepSeekMath-Base with a focus on mathematical instruction following. For each training instance, problems are concatenated at random until the combined sequence length approaches a 4,000 token limit. The model undergoes 500 training steps, using a batch size of 256 and a fixed learning rate of 5e-5.

Assessment of mathematical capability is performed with and without external tool access, using four distinct quantitative reasoning benchmarks in both English and Chinese. Comparative evaluations are drawn against state-of-the-art models available at the time, including:

\begin{itemize}[topsep=0pt]
    \item \textbf{Closed-source models} include: (1) the GPT model series, highlighting GPT-4 \citep{gpt4} and GPT-4 Code Interpreter~\footnote{\url{https://openai.com/blog/chatgpt-plugins\#\#code-interpreter}} for their advanced capabilities, (2) Gemini Ultra and Pro \citep{gemini}, (3) Inflection-2 \citep{inflection-2}, (4) Grok-1~\footnote{\url{https://x.ai/model-card}}, as well as (5) Baichuan-3~\footnote{\url{https://www.baichuan-ai.com}}, and (6) the recently released GLM-4~\footnote{\url{https://open.bigmodel.cn/dev/api\#glm-4}} from the GLM family \citep{glm}, all developed by leading entities, including notable Chinese technology companies.

These models are designed for broad applicability and have generally been subjected to multiple alignment processes.
    \item \textbf{Open-source models} encompass: general-purpose models such as (1) DeepSeek-LLM-Chat 67B \citep{deepseek-llm}, (2) Qwen 72B \citep{qwen}, (3) SeaLLM-v2 7B \citep{seallm}, and (4) ChatGLM3 6B \citep{chatglm3}. Additionally, there are models optimized for mathematical reasoning: (5) InternLM2-Math 20B~\footnote{\url{https://github.com/InternLM/InternLM-Math}}, an extension of InternLM2 with additional mathematics-specific pretraining and instruction tuning; (6) Math-Shepherd-Mistral 7B, which applies PPO training \citep{schulman2017proximal} to Mistral 7B \citep{mistral} using a process-supervised reward model; (7) the WizardMath series \citep{wizardmath}, which enhances the mathematical capabilities of Mistral 7B and Llama-2 70B \citep{llama2} via evolve-instruct—an instruction tuning strategy employing AI-generated instructions—and PPO on datasets primarily drawn from GSM8K and MATH; (8) MetaMath 70B \citep{metamath}, created by fine-tuning Llama-2 70B on an enriched GSM8K and MATH corpus; (9) ToRA 34B \cite{tora}, which involves fine-tuning CodeLlama 34B for tool-assisted mathematical reasoning; and (10) MAmmoTH 70B \citep{MathInstruct}, a version of Llama-2 70B tuned on MathInstruct through instruction tuning.
\end{itemize}

According to Table \ref{tab:sft_rl_math}, when model evaluation prohibits tool use, DeepSeekMath-Instruct 7B delivers strong performance on multi-step reasoning. In particular, on the competition-grade MATH benchmark, this model outperforms all open-source models as well as most proprietary systems (including Inflection-2 and Gemini Pro), exceeding them by a minimum margin of 9\% absolute—even when compared to considerably larger models like Qwen 72B or models specially reinforced for mathematics such as WizardMath-v1.1 7B. DeepSeekMath-Instruct’s performance on MATH matches that of leading Chinese proprietary models like GLM-4 and Baichuan-3, though it does not reach the levels achieved by GPT-4 or Gemini Ultra.

In the evaluation scenario that allows for a combination of natural language reasoning and program-based tool usage, DeepSeekMath-Instruct 7B attains nearly 60\% accuracy on MATH, setting a new high among open-source models. On additional benchmarks, its performance is comparable to that of DeepSeek-LLM-Chat 67B, the prior best-performing model in this category despite being ten times larger.
\section{Reinforcement Learning}

\subsection{Group Relative Policy Optimization} Reinforcement learning (RL) has shown substantial effectiveness in augmenting the mathematical reasoning capabilities of large language models (LLMs) after the Supervised Fine-Tuning (SFT) phase \citep{wang2023math,wizardmath}. We now introduce our proposed reinforcement learning algorithm—Group Relative Policy Optimization (GRPO)—which emphasizes efficiency and efficacy.

\subsubsection{From PPO to GRPO} Proximal Policy Optimization (PPO) \citep{schulman2017proximal} is a widely adopted actor-critic RL method in the RL fine-tuning phase for LLMs \citep{ouyang2022training}. Specifically, PPO optimizes LLMs by maximizing the following surrogate objective function:

\begin{equation}
\footnotesize
    \mathcal{J}_{PPO}(\theta) = \mathbb{E}{[q \sim P(Q), o \sim \pi_{\theta_{old}}(O|q)]} \frac{1}{|o|} \sum_{t=1}^{|o|} \min \left[ \frac{\pi_\theta(o_{t} | q, o_{<t})}{\pi_{\theta_{old}}(o_{t} | q, o_{<t})} A_{t}, \text{clip} \left( \frac{\pi_\theta(o_{t} | q, o_{<t})}{\pi_{\theta_{old}}(o_{t} | q, o_{<t})}, 1 - \varepsilon, 1 + \varepsilon \right)

Here, $\pi_{\theta}$ and $\pi_{\theta_{old}}$ represent the current and previous policy networks, respectively, while $q$ and $o$ denote questions and their associated outputs, sampled from the question dataset and according to the old policy $\pi_{\theta_{old}}$. The parameter $\varepsilon$ serves as a clipping hyper-parameter in PPO, contributing to the stability of training. The term $A_t$ indicates the advantage, estimated by Generalized Advantage Estimation (GAE) \citep{gae}, which is based on rewards $\{r_{\ge t}\}$ and an approximated value function $V_{\psi}$. Consequently, PPO requires learning a value function in tandem with the policy network. To guard against reward model overfitting, it is common to incorporate a per-token KL divergence penalty in the reward relative to a reference policy at each token \citep{ouyang2022training}, formalized as

\begin{equation}
    r_{t} = r_\varphi(q, o_{\le t}) - \beta \log\frac{\pi_{\theta}(o_{t}|q, o_{<t})}{\pi_{ref}(o_{t}|q, o_{<t})},
\label{eq:PPO-reward}
\end{equation}

where $r_\varphi$ represents the reward model, $\pi_{ref}$ denotes the reference policy (often initialized as the SFT model), and $\beta$ is the scaling factor for the KL regularization term.

Given that the value function employed in PPO typically matches the scale of the policy network, it imposes considerable computational and memory demands. Additionally, PPO regards the value function as a baseline for advantage computation to reduce variance, but in LLM applications, the reward model generally supplies a reward only for the last token, making it more challenging to train a value function that yields precise values at each token. In light of these issues, and as illustrated in Figure \ref{fig:grpo}, we introduce Group Relative Policy Optimization (GRPO). GRPO removes the requirement for learning a separate value function as in PPO by using the mean reward from a collection of sampled outputs—generated for the same question—as the baseline. Concretely, for each question $q$, a set of outputs $\{o_1, o_2, \cdots, o_G\}$ are sampled from $\pi_{\theta_{old}}$, and the policy is optimized to maximize the following target:

\begin{equation}
\footnotesize
\begin{split}
    \mathcal{J}_{GRPO}(\theta) &= \mathbb{E}{[q \sim P(Q), \{o_i\}_{i=1}^G \sim \pi_{\theta_{old}}(O|q)]}  \\
    & \frac{1}{G}\sum_{i=1}^G\frac{1}{|o_i|} \sum_{t=1}^{|o_i|} \left\{ \min \left[ \frac{\pi_\theta(o_{i,t} | q, o_{i,<t})}{\pi_{\theta_{old}}(o_{i,t} | q, o_{i,<t})} \hat{A}_{i,t}, \text{clip} \left( \frac{\pi_\theta(o_{i,t} | q, o_{i,<t})}{\pi_{\theta_{old}}(o_{i,t} | q, o_{i,<t})}, 1 - \varepsilon, 1 + \varepsilon \right)  \hat{A}_{i,t} \right] - \beta \mathbb{D}_{KL}\left[\pi_{\theta} || \pi_{ref}\right]\right\} ,
\end{split}
\label{eq:GRPO-obj}
\end{equation}

Here, $\varepsilon$ and $\beta$ continue to play the roles of tunable hyper-parameters, and $\hat{A}_{i,t}$ is the advantage computed based solely on the relative performance among outputs within each group—a procedure explained in detail in subsequent sections. GRPO’s group-based method for estimating advantages matches the comparative approach often used when training reward models, as these models are generally constructed from datasets containing paired output comparisons for the same prompt. Furthermore, in contrast to PPO, which embeds the KL penalty within the reward, GRPO directly penalizes the KL divergence between the candidate and reference policy in the loss itself, thereby simplifying the evaluation of $\hat{A}_{i,t}$. Importantly, the KL divergence term utilized here is not identical to the one found in (\ref{eq:PPO-reward}); instead, we compute it using the following unbiased estimator as presented in \

\begin{algorithm}[t]
  \small
  \caption{Iterative Group Relative Policy Optimization}
  \textbf{Input} initial policy model $\pi_{\theta_{\text{init}}}$; reward models $r_\varphi$; task prompts $\mathcal{D}$; 
  hyperparameters $\varepsilon$, $\beta$, $\mu$
  \begin{algorithmic}[1]
    \State Initialize the policy model: $\pi_\theta \leftarrow \pi_{\theta_{\text{init}}}$
    \For{iteration = 1, \dots, I}
       \State Set the reference model: $\pi_{ref} \leftarrow \pi_{\theta}$
      \For{step = 1, \dots, M}
      \State Draw a mini-batch $\mathcal{D}_b$ from $\mathcal{D}$
      \State Assign the current policy model to $\pi_{\theta_{old}} \leftarrow \pi_{\theta}$ 
      \State For each question $q \in \mathcal{D}_b$, sample $G$ outputs $\{o_i\}_{i=1}^G$ from $\pi_{\theta_{old}} (\cdot \mid q)$
      \State Evaluate all sampled outputs $\{o_i\}$ using $r_{\varphi}$ to obtain rewards $\{r_i\}_{i=1}^{G}$
      \State Estimate group-relative advantages $\hat{A}_{i,t}$ for the $t$-th token in each $o_i$
      \For{GRPO iteration = 1, \dots, $\mu$}
        \State Update $\pi_\theta$ by maximizing the GRPO objective as shown in Equation \ref{eq:GC-GRPO}
      \EndFor
    \EndFor \State Continuously train $r_\varphi$ using a replay-based update mechanism.
    \EndFor 
  \end{algorithmic}
  \textbf{Output} $\pi_\theta$
  \label{alg:iter-grpo}
\end{algorithm}

\subsubsection{Outcome Supervision RL with GRPO} 
Formally, for any question $q$, generate a collection of outputs $\{o_1, o_2, \cdots, o_G\}$ from the previous policy model $\pi_{\theta_{old}}$. These outputs are evaluated by the reward model, producing $G$ corresponding reward values $\mathbf{r}=\{r_1, r_2, \cdots, r_G\}$. The obtained rewards are standardized by removing the mean and scaling by the standard deviation of the group. In the outcome supervision setting, each output $o_i$ receives its normalized reward at completion, which is then assigned as the advantage $\hat{A}_{i, t}$ to every token in that output: $\hat{A}_{i, t} = \widetilde{r}_i = \frac{r_i- {\rm mean}(\mathbf{r})}{{\rm std}(\mathbf{r})}$. The policy model is subsequently improved by maximizing the objective given in equation (\ref{eq:GRPO-obj}).

\subsubsection{Process Supervision RL with GRPO} 
Outcome-based supervision only delivers a terminal reward per output, which can be limiting when guiding policy learning for tasks involving multiple reasoning steps. As in \cite{wang2023math}, we also utilize process supervision, where feedback is provided after each intermediate reasoning step. Specifically, for a prompt $q$ and $G$ sampled outputs $\{o_1, o_2, \cdots, o_G\}$, the process reward model scores each individual step, producing sets of step-wise rewards: $\mathbf{R} = \{ \{r_1^{index(1)},\cdots,r_1^{index(K_1)}\}, \cdots,  \{r_G^{index(1)},\cdots,r_G^{index(K_G)}\} \}$, with $index(j)$ denoting the terminal token index of step $j$, and $K_i$ the number of steps in output $i$. Each reward is normalized against the collective mean and standard deviation: $\widetilde{r}_i^{index(j)} = \frac{r_i^{index(j)} - {\rm mean(\mathbf{R})}}{{\rm std(\mathbf{R})}}$.
The advantage for each token is then defined as the sum of future normalized rewards over subsequent steps: $\hat{A}_{i, t} = \sum_{index(j) \ge t} \widetilde{r}_i^{index(j)}$. The policy is then optimized according to the objective in equation (\ref{eq:GR

\subsubsection{Iterative RL with GRPO} During the reinforcement learning process, the original reward model may eventually become inadequate for supervising the evolving policy model. To address this, we investigate an iterative reinforcement learning scheme leveraging GRPO. As outlined in Algorithm \ref{alg:iter-grpo}, iterative GRPO repeatedly regenerates training data for the reward model using samples drawn from the current policy model, retraining the existing reward model each time through a replay buffer that maintains 10\% of previous data. Afterward, the updated policy model is used as the new reference model, and further training of the policy model proceeds using the newly updated reward model.

\subsection{Training and Evaluating DeepSeekMath-RL}

Our reinforcement learning experiments are conducted using the DeepSeekMath-Instruct 7B architecture. The RL training utilizes chain-of-thought formatted questions drawn from the GSM8K and MATH subsets present in the SFT dataset, encompassing approximately 144,000 samples. We intentionally omit other SFT questions to study the effects of RL in benchmark scenarios where no related data is available throughout the RL process. Following the methodology in \citep{wang2023math}, we curate the reward model's training data. The initial reward model is trained on DeepSeekMath-Base 7B with a learning rate of 2e-5. In our GRPO framework, we specify the policy model's learning rate as 1e-6, and the KL penalty is set at 0.04. For each input question, 64 distinct outputs are sampled. The maximum sequence length for training is fixed at 1024 tokens, with a batch size also set to 1024. After every exploration iteration, the policy model undergoes only a single update. Evaluation of DeepSeekMath-RL 7B adheres to the same protocol as DeepSeekMath-Instruct 7B. Specifically, for DeepSeekMath-RL 7B, the GSM8K and MATH datasets featuring chain-of-thought reasoning are categorized as in-domain tasks, while all remaining benchmarks are considered out-of-domain tasks.

Table \ref{tab:sft_rl_math} presents the results achieved by both open-source and proprietary models employing chain-of-thought as well as tool-augmented reasoning on English and Chinese evaluation datasets. Our observations indicate: 1) DeepSeekMath-RL 7B achieves 88.2\% accuracy on GSM8K and 51.7\% on MATH, both with chain-of-thought reasoning. These results not only exceed all open-source models within the 7B to 70B parameter range but also outperform the majority of closed-source alternatives. 2) Notably, DeepSeekMath-RL 7B is fine-tuned solely on chain-of-thought formatted instruction data from GSM8K and MATH, building on DeepSeekMath-Instruct 7B. In spite of the limited diversity of training data, DeepSeekMath-RL 7B demonstrates improved performance over DeepSeekMath-Instruct 7B across every evaluation metric, underlining the advantage conferred by reinforcement learning.

\section{Discussion}
In this part, we discuss insights gained through our pre-training and RL experiments.

\subsection{Lessons Learnt in Pre-Training}

We begin by outlining our pre-training experiences. Unless explicitly stated otherwise, the training configurations follow those described in Section \ref{sec:quality-policy}. For this discussion, references to the DeepSeekMath Corpus pertain to the second iteration dataset, consisting of 89B tokens.

\subsubsection{Code Training Benefits Mathematical Reasoning}

There exists a widely held, though largely untested, hypothesis that exposure to code benefits reasoning capabilities. Our findings contribute evidence in support of this claim within the scope of mathematics: training with code data improves models' proficiency in mathematical reasoning, both in standard problem-solving and when leveraging tools.

To investigate the influence of code data on mathematical reasoning, we performed experiments comparing two-stage and single-stage training regimes:

\textbf{Two-Stage Training}

\begin{itemize}[topsep=0pt]
    \item \textbf{Code Training for 400B Tokens $\rightarrow$ Math Training for 150B Tokens}: The DeepSeek-LLM 1.3B model is subjected to pretraining on 400 billion code tokens, which is then followed by further training on 150 billion tokens sourced from mathematical data;
    \item \textbf{General Training for 400B Tokens $\rightarrow$ Math Training for 150B Tokens}: For comparison, an alternative setup uses general-domain tokens (drawn from a comprehensive general-purpose corpus curated by DeepSeek-AI) instead of code tokens in the initial training stage, with the goal of elucidating whether code pretraining provides a particular advantage over general text in enhancing mathematical reasoning capabilities.
\end{itemize}

\textbf{One-Stage Training}

\begin{itemize}[topsep=0pt]
    \item \textbf{Math Training for 150B Tokens}: Here, DeepSeek-LLM 1.3B undergoes training exclusively on 150 billion mathematical tokens;
    \item \textbf{Training on a mixture of 400B Code Tokens and 150B Math Tokens}: Given that subsequent mathematical training after code pretraining can adversely affect performance on coding, we explore the impact of training the model on a combined set of code and math tokens in a single stage to evaluate whether this strategy promotes mathematical reasoning and helps counteract catastrophic forgetting.
\end{itemize}

\paragraph{Results}
Table \ref{tab:code-math} and Table \ref{tab:code-math-general-eval} present the empirical results from various training protocols.

Pretraining on code tokens contributes positively to mathematical reasoning facilitated by program synthesis in both two-stage and one-stage regimes. The findings in Table \ref{tab:code-math} indicate that, in the two-stage approach, code pretraining alone markedly boosts the model’s proficiency at tackling GSM8K and MATH tasks using Python. A subsequent phase of math-specific training provides additional performance gains. Notably, the one-stage protocol that interleaves code and math tokens reduces the severity of catastrophic forgetting observed in sequential training and also leads to improvements in coding (see Table \ref{tab:code-math-general-eval}) as well as in program-supported mathematical reasoning (see Table \ref{tab:code-math}).

Code pretraining also enhances mathematical reasoning performance when the model does not rely on external tools. In the sequential training context, early exposure to code data yields moderate improvements and amplifies the benefits of later math training, culminating in optimal results. Conversely, merging code and math tokens during a single-stage process appears to diminish the model’s capability for mathematical reasoning unaided by tools. A plausible explanation is that the 1.3B parameter DeepSeek-LLM may not be large enough to efficiently encode and reconcile the combined complexities of both coding and mathematical knowledge within a single training run.

\subsubsection{ArXiv Papers Exhibit Limited Effectiveness for Enhancing Mathematical Reasoning}

Despite the widespread use of ArXiv papers as a significant source in the pre-training datasets for mathematical models \citep{minerva,gpt-f,llemma,mathpile}, their actual contribution to the advancement of mathematical reasoning abilities remains underexplored. In contrast to what might be expected, our findings indicate that incorporating ArXiv papers does not substantially enhance mathematical reasoning skills. We evaluated this by training models of varying scales—including DeepSeek-LLM 1.3B and DeepSeek-Coder-Base-v1.5 7B \citep{deepseek-coder}—using different versions of the ArXiv dataset that had undergone separate pre-processing workflows:

\begin{itemize}[topsep=0pt]
    \item \textbf{MathPile} \citep{mathpile}: an 8.9B-token dataset constructed via rigorous cleaning and filtering heuristics, with scientific ArXiv papers representing over 85\% of the content;
    \item \textbf{ArXiv-RedPajama} \citep{redpajama}: a corpus of all ArXiv LaTeX documents with extraneous preambles, comments, macros, and bibliographic entries stripped, totaling 28.0B tokens.
\end{itemize}

For our analysis, DeepSeek-LLM 1.3B was trained on each ArXiv corpus for 150B tokens, while DeepSeek-Coder-Base-v1.5 7B was trained for 40B tokens. Results reveal that using ArXiv papers alone does not bolster mathematical reasoning; the models did not exhibit notable improvements and, in some instances, their performance declined across a range of mathematical tasks of diverse complexity levels considered in this work. The evaluated tasks encompass quantitative reasoning problems such as GSM8K and MATH (see Table \ref{tab:arxiv-cot}), multiple-choice assessments like MMLU-STEM (Table \ref{tab:arxiv-cot}), and formal mathematics problems such as miniF2F (refer to Table \ref{tab:arxiv_atp}).

Nevertheless, these observations must be interpreted with caution given certain constraints of our study. Specifically, we have not yet investigated:

\begin{itemize}[topsep=0pt]
    \item How ArXiv-sourced tokens affect mathematical tasks outside the scope of this work, for instance, the task of theorem informalization, which involves rephrasing formal statements or proofs in a more informal manner;
    \item The interaction between ArXiv tokens and other varieties of training data;
    \item Whether the utility of ArXiv data emerges more prominently at increased model sizes.
\end{itemize}

Consequently, these questions highlight areas for continued investigation in future research.

\subsection{Insights of Reinforcement Learning}
\subsubsection{Towards to a Unified Paradigm} Here, we introduce a unified framework for analyzing multiple training strategies such as SFT, RFT, DPO, PPO, and GRPO, and subsequently run experiments to scrutinize key aspects within this framework. In general, the gradient of a given training method with respect to the parameter $\theta$ can be formalized as follows:

\begin{equation}
    \nabla_{\theta}\mathcal{J}_{\textcolor{red}{\mathcal{A}}}(\theta) = \mathbb{E}[\underbrace{(q,o) \sim \textcolor{red}{\mathcal{D}}}_{Data \ Source}]\left( \frac{1}{|o|} \sum_{t=1}^{|o|}  \underbrace{GC_{{\mathcal{A}}}(q, o, t, \textcolor{red}{\pi_{{rf}}})}_{Gradient \ Coefficient}  \nabla_{\theta}\log \pi_{\theta}(o_t | q, o_{<t})\right).
\label{eq:objective}
\end{equation}

This formulation comprises three fundamental elements: (1) \textit{Data Source $\mathcal{D}$}, specifying the set of samples used during training; (2) \textit{Reward Function $\pi_{{rf}}$}, which delivers the evaluative feedback that drives model improvement; and (3) \textit{Algorithm $\mathcal{A}$}, which translates the data and rewards into the gradient coefficient $GC$, effectively shaping the learning dynamics via the applied gradients. We categorize and examine various widely used approaches according to this unifying framework:

\begin{itemize}
    \item  \textbf{Supervised Fine-tuning (SFT)}: SFT adapts a pretrained model through training on datasets curated by human annotators.
    \item \textbf{Rejection Sampling Fine-tuning (RFT)}: RFT continues the fine-tuning of the SFT-initialized model by leveraging outputs generated from SFT questions; these outputs are filtered to retain only those with correct answers before being used for further training.
    \item \textbf{Direct Preference Optimization (DPO)}: DPO further enhances the SFT-trained model by conducting fine-tuning on an expanded set of model-generated responses, with learning guided by pairwise DPO loss.
    \item \textbf{Online Rejection Sampling Fine-tuning (Online RFT)}: Unlike traditional RFT, Online RFT continually updates the policy model initialized from SFT by training on fresh, model-generated outputs in real time.
    \item \textbf{PPO/GRPO}: With PPO/GRPO, the policy is initialized using the SFT model and strengthened through outputs sampled on-the-fly from the evolving policy itself.
\end{itemize}

The elements of these approaches are outlined in Table \ref{tab:data_policy}. For an in-depth explanation of the derivation, see Appendix \ref{app:analysis-rl}.

\paragraph{Remarks on Data Source} We categorize data sources as either online or offline sampling. Online sampling refers to collecting training data from the live outputs of the current policy model as it explores, while offline sampling utilizes data generated by the original SFT model before policy training commences. RFT and DPO are both based on offline data acquisition, whereas Online RFT and GRPO adopt online data collection strategies.

From Figure \ref{fig:rl-analysis}, it is evident that Online RFT considerably surpasses RFT on both benchmarks. Initially, Online RFT and RFT display similar performance levels, but as training progresses, Online RFT establishes a clear lead, highlighting the efficacy of online training. This observation aligns with expectations: early in training, there is little distinction between the actor and the SFT model, so their respective sampled data also differ minimally. As training advances, data generated by the evolving actor diverges more notably, and continuously updated data from online sampling provides pronounced benefits.

\paragraph{Remarks on Gradient Coefficient} The algorithm translates reward signals into gradient coefficients to update model parameters. Our experiments distinguish the reward function into `Rule'—which assesses response quality by answer correctness—and `Model,' where a learned reward model assigns scores to each response. The reward model is trained using labels generated via rule-based evaluation. As reflected in Equations \ref{eq:GC-RFT} and \ref{eq:GC-GRPO}, a fundamental difference exists between GRPO and Online RFT: GRPO uniquely adapts its gradient coefficient according to the reward output by the reward model. This mechanism enables nuanced positive and negative updates, reflecting varying degrees of reward. Online RFT, conversely, lacks this granularity; incorrect responses receive no explicit penalty, and all correct responses are rewarded with the same intensity.

As illustrated in Figure \ref{fig:rl-analysis}, GRPO outperforms online RFT, emphasizing the advantage of modulating the coefficients for positive and negative gradients. Moreover, GRPO+PS achieves better results than GRPO+OS, which underscores the effectiveness of employing step-sensitive, fine-grained gradient coefficients. Additionally, we investigate iterative RL by performing two cycles of training in our experiments. According to Figure \ref{fig:iter-rl}, the iterative application of RL markedly enhances performance, with the most substantial improvement occurring during the initial iteration.

\subsubsection{Why RL Works?} In this study, reinforcement learning is applied to a selected subset of instruction tuning data, leading to considerable gains over the instruction-tuned model baseline. To gain insight into the mechanisms behind RL’s success, we assess both Pass@K and Maj@K accuracy for the Instruct and RL-trained models across two evaluation sets. As depicted in Figure \ref{fig:maj-pass}, RL training yields notable improvements in Maj@K while leaving Pass@K mostly unaffected. These results suggest that the gains from RL stem from making the output distribution more reliable—\textbf{suggesting that progress is achieved by promoting a correct answer among the TopK predictions, rather than by strengthening fundamental abilities.} Likewise, prior work \citep{wang2023making} reported a \textbf{misalignment problem} with reasoning in SFT models, showing that preference-based alignment strategies can improve reasoning accuracy \citep{yuan2023rrhf,song2023preference,wang2023making}.

\subsubsection{How to Achieve More Effective RL?}
\label{sec:effec_rl}
Our findings confirm the strong performance of RL approaches in mathematical reasoning. We further introduce a unified conceptual framework encompassing various popular training strategies, interpreting them as either explicit or reduced forms of RL. Equation \ref{eq:objective} formalizes this by distinguishing three principal aspects: Data Source, Algorithm, and Reward Function. We discuss several possible avenues for advancing each of these components.

\paragraph{Data Source} The choice of data is foundational to all training paradigms. In RL, we specifically define the data source as the collection of unlabeled questions along with responses sampled from the policy network. Our experiments restrict the source to queries from instruction tuning and apply basic nucleus sampling to generate outputs. We hypothesize that this limitation explains why improvements in our RL approach are mainly restricted to Maj@K. Looking ahead, we intend to extend our RL pipeline to out-of-distribution question prompts and to leverage \textbf{more sophisticated sampling (decoding) approaches}, such as those employing tree-search \citep{yao2023tree}. In addition, developments in \textbf{efficient inference strategies} \citep{xia-etal-2023-speculative,leviathan2023fast,kwon2023efficient,xia2024unlocking}—which impact how effectively policy models can explore—are also crucial to the advancement of RL methods.

\paragraph{Algorithms} Algorithms utilize both input data and reward signals to compute gradient coefficients, thereby updating model parameters. In line with Equation \ref{eq:objective}, contemporary approaches, to varying degrees, fundamentally \textbf{TRUST} the reward function signal for modulating the conditional probability of a given token—either increasing or decreasing it. Nevertheless, the infallibility of the reward signal cannot be guaranteed, particularly in highly intricate scenarios. For instance, the PRM800K datasets \citep{lightman2023let}, despite being rigorously labeled by expert annotators, still exhibit a roughly 20\% error rate in annotations\footnote{\url{https://github.com/openai/prm800k/issues/12\#issuecomment-1728491852}}. This observation motivates our investigation into reinforcement learning algorithms designed to be resilient in the face of noisy reward signals. We contend that such \textbf{WEAK-TO-STRONG} \citep{burns2023weak} alignment frameworks could fundamentally reshape the underlying learning paradigms.

\paragraph{Reward Function} The reward function serves as the principal mechanism for delivering learning signals during training. In reinforcement learning, this typically takes the form of a neural reward model. We identify three critical avenues for advancement in reward modeling: 1) \textbf{Improving the generalization capacity of the reward model.} It is imperative for the reward model to effectively generalize to out-of-distribution problems and sophisticated decoded responses; failing this, reinforcement learning risks merely consolidating existing LLM distributions rather than achieving substantive progress; 2) \textbf{Capturing the uncertainty in the reward model.} Accurately characterizing uncertainty may facilitate integration between underpowered reward models and weak-to-strong alignment approaches; 3) \textbf{Constructing process-level reward models of high fidelity with greater efficiency} so that they can supply granular training signals during reasoning tasks \citep{lightman2023let,wang2023math}.

\section{Conclusion, Limitation, and Future Work}

We introduce DeepSeekMath, which surpasses all existing open-source models on the competitive MATH benchmark and achieves results that are close to those of proprietary systems. DeepSeekMath~is built upon DeepSeek-Coder-v1.5 7B, undergoing continuous training for 500B tokens—of which a significant portion, 120B tokens, are mathematics-focused data extracted from Common Crawl. Through comprehensive ablation experiments, we demonstrate that online sources contribute rich, high-quality mathematical datasets, whereas arXiv does not appear to provide as much value as initially anticipated. We also propose Group Relative Policy Optimization (GRPO), an adaptation of Proximal Policy Optimization (PPO), which enhances mathematical reasoning performance while reducing memory usage. Experimental analysis indicates that GRPO can still yield notable gains even when DeepSeekMath-Instruct 7B is already performing at a high level on standard benchmarks. Furthermore, we provide a cohesive framework for interpreting various approaches and outline promising avenues for further advancing reinforcement learning efficiency.

Despite its strong results on benchmarks requiring quantitative reasoning, DeepSeekMath~still exhibits weaker proficiency in areas such as geometry and theorem proving compared to closed-source alternatives. For example, in preliminary testing, the model struggled with geometric problems involving triangles and ellipses, suggesting the possibility of bias in the chosen training and fine-tuning datasets. Moreover, due to the model’s limited size, DeepSeekMath~underperforms relative to GPT-4 on few-shot learning tasks: while GPT-4 benefits noticeably from few-shot contexts, DeepSeekMath~does not show substantial improvement when moving from zero-shot to few-shot settings. Moving forward, we plan to refine our engineered data selection methodology in order to construct an even higher-quality pre-training corpus. Additionally, we will investigate the directions highlighted in Section \ref{sec:effec_rl} to further enhance reinforcement learning strategies for large language models.